The questions in this section build a two-player tic-tac-toe web app using Lustre, the same framework as the rock-paper-scissors app from lab. They do not depend on the lab.

\subsection*{The Model and Messages}

In a Lustre app, the \emph{model} holds the entire state of the app, \texttt{view} turns the model into a tree of elements to display, and \texttt{update} takes the model and a \emph{message} and returns a new model. When the user clicks a button, Lustre sends that button's message to \texttt{update}, then calls \texttt{view} on the new model and changes the page to match.

Here, the model is a board of 9 cells, numbered \texttt{0} to \texttt{8} from left to right, top to bottom. Clicking cell \texttt{i} sends \lstinline{UserClickedCell(i)}, and clicking ``New game'' sends \lstinline{UserClickedNewGame}.
\begin{lstlisting}
pub type Player { X O }
pub type Cell { Empty Taken(Player) }
pub type Model { Model(board: List(Cell)) }
pub type Msg { UserClickedCell(Int) UserClickedNewGame }
\end{lstlisting}
For example, after \texttt{X} takes the center and \texttt{O} takes the top-left corner, the board is \lstinline{[Taken(O), Empty, Empty, Empty, Taken(X), Empty, Empty, Empty, Empty]}. \texttt{X} always moves first, and there is no winner: the game continues until the board is full.

\subsection*{Provided Code}

Each Lustre element is made by a function that takes a list of \emph{attributes} and a list of \emph{children}. For example, \lstinline{html.button([attribute.class("cell"), event.on_click(UserClickedNewGame)], [html.text("X")])} is a button with class \texttt{"cell"} and text \texttt{X} that sends \lstinline{UserClickedNewGame} when clicked. \lstinline{attribute.disabled(True)} makes a button unclickable.
\begin{itemize}
  \item \lstinline{init(Nil)} returns an empty board, and \lstinline{cell_at(board, i)} returns cell \texttt{i} of \texttt{board}.
  \item \texttt{view} shows the status, the board, and a ``New game'' button. \texttt{view\_status} shows whose turn it is, such as \texttt{"X to move"}, or \texttt{"Board full"}.
  \item \texttt{view\_board} puts one button per cell in a \lstinline{div} with class \texttt{"board"}, which CSS displays as a 3-by-3 grid. \lstinline{list.index_map} is like \lstinline{list.map}, but calls \lstinline{view_cell(cell, i)} with each cell and its index.
\end{itemize}
\begin{lstlisting}
pub fn init(_flags) -> Model {
  Model(board: list.repeat(Empty, 9))
}

pub fn cell_at(board: List(Cell), i: Int) -> Cell {
  case list.drop(board, i) {
    [cell, ..] -> cell
    [] -> Empty
  }
}

pub fn player_name(player: Player) -> String {
  case player {
    X -> "X"
    O -> "O"
  }
}

pub fn view(model: Model) -> Element(Msg) {
  html.div([attribute.class("game")], [
    view_status(model),
    view_board(model),
    html.button([event.on_click(UserClickedNewGame)], [html.text("New game")]),
  ])
}

pub fn view_status(model: Model) -> Element(Msg) {
  let message = case list.contains(model.board, Empty) {
    True -> player_name(current_player(model.board)) <> " to move"
    False -> "Board full"
  }
  html.p([attribute.class("status")], [html.text(message)])
}

pub fn view_board(model: Model) -> Element(Msg) {
  html.div([attribute.class("board")], list.index_map(model.board, view_cell))
}
\end{lstlisting}